\newpage
\section{Solving Systems of Linear Equations}
Consider a system of linear equations with $n$ unknowns and $m$ equations:
$$(*) = 
\begin{cases}
\begin{aligned}
a_{11}x_1 & + a_{12}x_2 + \cdots + a_{1n}x_n && = b_1 \\
a_{21}x_1 & + a_{22}x_2 + \cdots + a_{2n}x_n && = b_2 \\
&\ \vdots && \ \vdots \\
a_{m1}x_1 & + a_{m2}x_2 + \cdots + a_{mn}x_n && = b_m
\end{aligned}
\end{cases}
$$

\begin{theorem}
For $(*)$ exactly one of the following holds:
\begin{itemize}
    \item $(*)$ has exactly one solution (consistent and independent)
    \item $(*)$ has infinitely many solutions (consistent and dependent)
    \item $(*)$ has no solutions (inconsistent)
\end{itemize}
\end{theorem}

\begin{example}
We could rewrite $(*)$ into its matrix form, $Ax = b$, where $A$ is the coefficient matrix and $b$ is the vector of constants,
$$\underbrace{\bmat{a_{11} & a_{12} & \cdots & a_{1n} \\ a_{21} & a_{22} & \cdots & a_{2n} \\ \vdots & \vdots & \ddots & \vdots \\ a_{m1} & a_{m2} & \cdots & a_{mn}}}_A \underbrace{\bmat{x_1 \\ x_2 \\ \vdots \\ x_n}}_x = \underbrace{\bmat{b_1 \\ b_2 \\ \vdots \\ b_n}}_b$$
and we can denote the augmented matrix $[A \mid b]$ as
$$\begin{amatrix}{4}
a_{11} & a_{12} & \cdots & a_{1n} & b_1 \\
a_{21} & a_{22} & \cdots & a_{2n} & b_2 \\
\vdots & \vdots & \ddots & \vdots & \vdots \\
a_{m1} & a_{m2} & \cdots & a_{mn} & b_n \\
\end{amatrix}
$$
Notice that we have placed the coefficients in the left side and is separated by a vertical bar.\footnote{This serves as a visual representation on what we're separating.} On the right side we have the constants.
\end{example}

\begin{example}
Consider the following system of equations
$$\begin{cases}
    w + 3x + y + z = 3 \\
    2w + 2x + y + 2z = 8 \\
    3w + x + 2y - z = -1
\end{cases}$$
\end{example}

We can rewrite this in matrix form as:
$$\left[ 
\begin{matrix}
1 & 3 & 1 & 1 \\
2 & 2 & 1 & 2 \\
3 & 1 & 2 & -1
\end{matrix}
\right]
\left[
\begin{matrix}
w \\
x \\
y \\
z \\
\end{matrix}
\right] = 
\left[
\begin{matrix}
3 \\
8 \\
-1
\end{matrix}
\right]
$$

Shown below is its form in augmented matrix:
$$
\begin{amatrix}{4}
1 & 3 & 1 & 1 & 3 \\
3 & 1 & 2 & -1 & -1 \\
\end{amatrix}$$

\newpage
\underline{\textbf{Recall}}

In your previous math classes (Pre-calculus), you have learned solving system of equations through elimination by addition. To solve a system of equations, the goal is to convert it to an easier form that is easy to solve.

\underline{\textbf{Question}.} Which processes will yield equivalent system of equations?

Suppose we have the following system of equations:
$$\begin{cases}
    x - y = -2 \\
    2x + 6y = 12
\end{cases}$$

I. Switching the position of the equations:
\begin{align*}
\begin{cases}
    x - y = -2 \\
    2x + 6y = 12
\end{cases} \longleftrightarrow\,\,\,\,\,
\begin{cases}
    2x + 6y = 12 \\
    x - y = -2
\end{cases}
\end{align*}

II. Multiplying by a nonzero constant,
\begin{align*}
\begin{cases}
    x - y = -2 \\
    2x + 6y = 12
\end{cases} \longleftrightarrow \,\,\,\,\,
\begin{cases}
    2x + 6y = 12 \\
    x - y = -2
\end{cases}
\end{align*}


III. Add a constant multiple of one equation to another
\begin{align*}
\begin{cases}
    x - y = -2 \\
    2x + 6y = 12
\end{cases} 
&\longleftrightarrow \,\,\,\,\,
\begin{cases}
    x - y = -2 \\
    8y = 16
\end{cases} \\[1em]
&\longleftrightarrow \,\,\,\,\,
\begin{cases}
    x - y = -2 \\
    y = 2
\end{cases}
\end{align*}
All of these operations will yield the equivalent system of equations!

Suppose we take the system of equations from above and convert it to an augmented matrix,
$$\begin{amatrix}{2}
1 & 1 & - 2\\ 0 & 1 & 2 \\
\end{amatrix}
$$
Notice that it has a staircase form shown, as 
$$
\begin{bNiceMatrix}
1 & 1 & -2 \\ 
0 & 1 & \hspace{0.3cm}2 
\CodeAfter
  \tikz \draw[red, semithick, fill=red, fill opacity=0.1] 
    (1-1.north west) -- (1-3.north east) -- 
    (1-3.south east) -- (2-3.south east) -- 
    (2-2.south west) -- (1-2.south west) -- 
    (1-1.south west) -- cycle;
\end{bNiceMatrix}
$$
This form is also known as the \href{https://math.stackexchange.com/questions/1628192/what-does-echelon-mean/1629207#1629207}{row echelon form}.
\vspace{1em}

\begin{definition}
An $m \times n$ matrix is said to be in \textbf{row echelon form} if the following holds:
\begin{itemize}
    \item All zero rows of $A$ (if any) are at the bottom of $A$;
    \item The leftmost nonzero entry of every nonzero row of $A$ is 1;
    \item The leading entry of a higher row is strictly to the left of the leading entry of the next row.
\end{itemize}
\end{definition}

\begin{remark}
$A \in M_{m \times n}(\mathbb{R)}$ is in RREF (reduced row echelon form) if $A$ is in REF (reduced echelon form) and every leading entry of $A$ is the only nonzero entry of its column.
\end{remark}

\newpage
\begin{remark}
Each equation in a system corresponds to a row in its augmented matrix. The following processes will now yield equivalent systems:
\begin{itemize}
    \item \textbf{Type I} - Switch two rows in augmented matrix;
    \item \textbf{Type II} - Multiply a row in augmented matrix by nonzero constant;
    \item \textbf{Type III} - Add a constant multiple of one row in the augmented matrix to another row.
\end{itemize}
\end{remark}

\subsection{Gauss-Jordan Reduction}
To solve the system $Ax = b$:
\begin{itemize}
    \item Form the augmented matrix $[A \mid b]$
    \item Use elementary row operations to bring augmented matrix to RREF
    \item Express the resulting matrix as a system of linear equations
    \item Use back-substitution
\end{itemize}

\begin{remark}
\textbf{Gaussian elimination} is the term when dealing with matrices in REF form.
\end{remark}

\begin{example}
Use Gauss-Jordan Reduction to find the solution set of the following equations:
$$\begin{cases}
    2x + 3y + z = 2 \\
    2x + 4y + 2z = 8 \\
    x - y - z = 4
\end{cases}$$
\end{example}

\textit{Solution.} 

Form the augmented matrix and apply elementary row operations,
\begin{align*}
    \begin{amatrix}{3}
    2 & 3 & 1 & 2 \\ 
    2 & 4 & 2 & 8 \\
    1 & -1 & -1 & 4
\end{amatrix} &\xrightarrow{R_1 \leftrightarrow R_3}
\begin{amatrix}{3}
    1 & -1 & -1 & 4 \\
    2 & 4 & 2 & 8 \\
    2 & 3 & 1 & 2
\end{amatrix} \\
\xrightarrow{\substack{R_2 \leftarrow R_2 - R_1 \\ R_3 \leftarrow R_3 - R_1}}
\begin{amatrix}{3}
    1 & -1 & -1 & 4 \\ 
    0 & 6 & 4 & 0 \\
    0 & 5 & 3 & -6
\end{amatrix}
& \xrightarrow{R_2 \leftarrow R_2 - R_3}
\begin{amatrix}{3}
    1 & -1 & -1 & 4 \\
    0 & 1 & 1 & 6 \\
    0 & 5 & 3 & -6
\end{amatrix} \\
\xrightarrow{R_3 \leftarrow R_3 - 5R_2}
\begin{amatrix}{3}
    1 & -1 & -1 & 4 \\
    0 & 1 & 1 & 6 \\
    0 & 0 & -2 & -36
\end{amatrix} 
& \xrightarrow{R_3 \leftarrow -\frac{1}{2}R_3}
\begin{amatrix}{3}
    1 & -1 & -1 & 4 \\
    0 & 1 & 1 & 6 \\
    0 & 0 & 1 & 18
\end{amatrix} \\
\xrightarrow{R_1 \leftarrow R_1 + R_2}
\begin{amatrix}{3}
    1 & 0 & 0 & 10 \\
    0 & 1 & 1 & 6 \\
    0 & 0 & 1 & 18
\end{amatrix} 
& \xrightarrow{R_2 \leftarrow R_2 - R_3}
\begin{amatrix}{3}
    1 & 0 & 0 & 10 \\
    0 & 1 & 0 & -12 \\
    0 & 0 & 1 & 18
\end{amatrix} \\
\Longrightarrow &\begin{cases}
    x = 10 \\ 
    y = -12 \\
    z = 18
\end{cases}
\end{align*}

Therefore, the solution set is $\boxed{\{(10, -12, 18)\}}$.

\newpage
\begin{example}
Use Gauss-Jordan Reduction to find the solution set of the following equations:
$$\begin{cases}
    w + x + 2y + 3z = 13 \\
    w - 2x + y + z = 8 \\
    3w + x + y - z = 1
\end{cases}$$
\end{example}

\textit{Solution.} 

Form the augmented matrix and apply elementary row operations,
\begin{align*}
\begin{amatrix}{4}
    1 & 1 & 2 & 3 & 13 \\
    1 & -2 & 1 & 1 & 8 \\
    3 & 1 & 1 & -1 & 1
\end{amatrix} &\xrightarrow{\substack{R_2 \leftarrow R_2 - R_1 \\ R_3 \leftarrow R_3 - 3R_1}} 
\begin{amatrix}{4}
    1 & 1 & 2 & 3 & 13 \\
    0 & -3 & -1 & -2 & -5 \\
    0 & -2 & -5 & -10 & -38 
\end{amatrix} \\
\xrightarrow{R_2 \leftarrow R_2 - R_3} 
\begin{amatrix}{4}
    1 & 1 & 2 & 3 & 13 \\
    0 & -1 & 4 & 8 & 33 \\0 & -2 & -5 & -10 & -38 
\end{amatrix} & \xrightarrow{R_3 \leftarrow R_3 - 2R_2}
\begin{amatrix}{4}
    1 & 1 & 2 & 3 & 13 \\
    0 & -1 & 4 & 8 & 33 \\
    0 & 0 & -13 & -26 & -104
\end{amatrix} \\
\xrightarrow{\substack{R_2 \leftarrow -R_2 \\ R_3 \leftarrow -\frac{1}{13}R_3}} 
\begin{amatrix}{4}
    1 & 1 & 2 & 3 & 13 \\
    0 & 1 & -4 & -8 & -33 \\
    0 & 0 & 1 & 2 & 8
\end{amatrix} & \xrightarrow{R_1 \leftarrow R_1 - R_2} 
\begin{amatrix}{4}
    1 & 0 & 6 & 11 & 46 \\
    0 & 1 & -4 & -8 & -33 \\
    0 & 0 & 1 & 2 & 8
\end{amatrix} \\
\xrightarrow{\substack{R_1 \leftarrow R_1 - 6R_3 \\ R_2 \leftarrow R_2 + 4R_3}}
\begin{amatrix}{4}
    1 & 0 & 0 & -1 & -2 \\
    0 & 1 & 0 & 0 & -1 \\
    0 & 0 & 1 & 2 & 8
\end{amatrix} & \Longrightarrow 
\begin{cases}
    w - z = -2 \\
    x = -1 \\
    y + 2z = 8
\end{cases}
\end{align*}
Let $z = t$, then, $w = t - 2, y = 8-2t$. Therefore, the solution set is $\boxed{\{(t - 2, -1, -2t + 8, t) \mid t \in \RR\}}$.

\begin{remark}
\leavevmode
\begin{enumerate}
    \item The variable $z$ in the system of equations above is called a \textbf{free variable}.
    \item A system $Ax = b$ having at least one solution has $k$ free variables, where $k$ is the number of columns in the RREF of $[A \mid b]$ having no leading entries. The system is said to have a $k$-parameter solution.
\end{enumerate}
\end{remark}

For the rest of the examples, I will not show the whole process.\footnote{Since this is TOO tedious to format.} It is up to the reader to verify that the solution is correct.

\begin{example}
Use Gaussian Elimination to find the solution set.
$$\begin{cases}
    x - y + 2z = 4 \\
    x + z = 6 \\
    2x - 3y + 5z = 4 \\
    3x + 2y - z = 1
\end{cases}$$

\textit{Solution.}

Form the augmented matrix and apply elementary row operations,
\begin{align*}
    \begin{amatrix}{3}
        1 & -1 & 2 & 4 \\
        1 & 0 & 1 & 6 \\
        2 & -3 & 5 & 4 \\
        3 & 2 & -1 & 1
    \end{amatrix} &\xrightarrow{\text{REF}}
    \begin{amatrix}{3}
        1 & -1 & 2 & 4 \\
        0 & 1 & -1 & 2 \\
        0 & 0 & 1 & 21/2 \\
        0 & 0 & 0 & 1
    \end{amatrix} \\
\end{align*}
$$\Longrightarrow \begin{cases}
        x - y + 2z = 4 \\
        y - z = 2 \\
        z = 21/2 \\
        0 = 1 \Longrightarrow \text{This is a contradiction!}
    \end{cases}$$
Since the solution is inconsistent, therefore, the solution set is $\emptyset$ or $\{\}$.
\end{example}

\begin{remark}
The system $Ax = b$ has no solution iff $\bmat{0 & 0 & \cdots & 0 & \mid & 1}$ is a row in an REF of $[A \mid b]$.
\end{remark}

\begin{example}
Solve for the solution of the following system of equations.
$$\begin{cases}
    2x_1 + x_2 - 4x_3 - 5x_4 + 5x_5 = 5 \\
    3x_1 - x_2 - 6x_3 - 5x_4 = 5 \\
    5x_1 - 2x_2 - 10x_3 - 8x_4 - x_5 = 8
\end{cases}$$

\textit{Solution.} 

Form the augmented matrix and apply elementary row operations,
$$\begin{amatrix}{5}
    2 & 1 & -4 & -5 & 5 & 5 \\
    3 & -1 & -6 & -5 & 0 & 5 \\
    5 & -2 & -10 & -8 & -1 & 8
\end{amatrix} \xrightarrow{\text{RREF}}
\begin{amatrix}{5}
    1 & -2 & -2 & 0 & -5 & 0 \\
    0 & 1 & 0 & -1 & 3 & 1 \\
    0 & 0 & 0 & 0 & 0 & 0\\
\end{amatrix}
$$
$$\Longrightarrow \begin{cases}
    x_1 - 2x_2 - 2x_3 - 5x_5 = 0 \\
    x_2 - x_4 + 3x_5 = 1
\end{cases}$$
Let $x_3 = r, x_4 = s, x_5 = t$, then, $x_2 = s - 3t + 1$ and $x_1 = 2r + 2s - t + 2$. Therefore, the solution set is $\boxed{\{(2r + 2s - t + 2, s - 3t + 1, r, s,t) \mid r, s, t \in \RR\}}$.
\end{example}

\newpage
\subsection{Exercises}

\begin{problem}
Use the Gauss-Jordan reduction or Gaussian elimination to solve the system of equations
$$\begin{cases}
    3x + 4y - z  = 8 \\
    3x + 8y - 2z = 8 \\
    6x + 8y - 2z = 3
\end{cases}$$
\end{problem}

\begin{problem}
Let $a \in \mathbb{R}$, and consider the linear system in the unknowns $x_1, x_2, x_3, x_4$:
$$\begin{cases}
    x_1 +2x_2 -x_3 = 1 \\ 
    2x_1 +4x_2 -x_3 + x_4 = 1 \\
    x_1 +2x_2 +x_3 +2x_4 = \alpha
\end{cases}$$

(a) Find the value of $\alpha$ that would make the system consistent.

% \textit{Answer.} $\alpha = -1$

(b) When the system is consistent, find its solution set.

% \textit{Answer.} Solution Set: $\{(-2s - t, s, -t - 1, t) \mid s, t \in \mathbb{R}\}$
\end{problem}

\begin{problem}
Let
$$A = \bmat{\frac{1}{4}\sqrt{3} + \frac{1}{2} & \frac{1}{4}\sqrt{3} - \frac{1}{2} & -\frac{1}{4}\sqrt{2} \\ \frac{1}{4}\sqrt{3} - \frac{1}{2} & \frac{1}{4}\sqrt{3} - \frac{1}{2} & -\frac{1}{4}\sqrt{2} \\ \frac{1}{4}\sqrt{2} & \frac{1}{4}\sqrt{2} & \frac{1}{2}\sqrt{3}}$$
Find all $3 \times 1$ matrices $x$ satisfying $Ax = x$. [\textit{Hint: The matrix equation is equivalent to $(A - I_3)x = 0_{3 \times 1}$.}]
\end{problem}

\begin{problem}
Solve for the following system of nonlinear equations for $x, y$, and $z$.
\begin{align*}
x^2 + y^2 + z^2 = 6 \\
x^2 - y^2 + 2z^2 = 2 \\
2x^2 + y^2 - z^2 = 3
\end{align*}
[\textit{Hint: Let $X = x^2, Y = y^2, Z = z^2$.}]
\end{problem}

\begin{problem}
Solve the following system of nonlinear equations for the unknown angles $\alpha, \beta, \gamma$, where 
$0 \leq \alpha \leq 2\pi, 0 \leq \beta \leq 2\pi$, and $0 \leq \gamma \leq \pi$.
\begin{align*}
2\sin\alpha - \cos\beta + 3\tan\gamma = 3 \\
4\sin\alpha + 2\cos\beta - 2\tan\gamma = 2 \\
6\sin\alpha - 3\cos\beta + \tan\gamma = 9
\end{align*}
\end{problem}